\documentclass[conference]{IEEEtran}
\usepackage{cite}
\usepackage{amsmath,amssymb,amsfonts}
\usepackage{algorithmic}
\usepackage{graphicx}
\usepackage{textcomp}
\usepackage{xcolor}
\usepackage{booktabs}
\usepackage{multirow}

\begin{document}

\title{Network-Based Analysis of Polymer-Clay Nanocomposite Mechanical Properties: A Modulus-Driven Clustering Approach}

\author{\IEEEauthorblockN{Author Name}
\IEEEauthorblockA{\textit{Department of Engineering} \\
\textit{University Name}\\
City, Country \\
email@university.edu}
}

\maketitle

\begin{abstract}
This study presents a novel graph-theoretic approach for analyzing the mechanical property enhancement patterns in polymer-clay nanocomposites. A comprehensive dataset of 927 nanocomposite samples was systematically organized into four clusters based on neat polymer matrix elastic modulus, ranging from soft elastomeric (0-0.1 GPa) to rigid polymer systems ($>$1.0 GPa). Bipartite graphs were constructed for each cluster to represent relationships between polymer matrices and their corresponding nanocomposite formulations, with edge weights quantifying elastic modulus improvements. Centrality metrics including degree, weighted degree, and betweenness centrality were computed to identify ``knowledge hubs'' representing polymer systems with substantial cumulative property improvements. Results reveal an inverse relationship between matrix stiffness and improvement magnitude, with soft elastomeric systems (C1-C2) achieving 288-351\% average improvements compared to 33-46\% for rigid polymers (C3-C4). Surface modification effectiveness was found to be stiffness-dependent, showing significant benefits in semi-soft thermoplastics (222.6\% higher improvement for modified systems in C2) but diminishing returns in rigid polymers. This network-based methodology provides a systematic framework for guiding material selection in nanocomposite design.
\end{abstract}

\begin{IEEEkeywords}
polymer nanocomposites, clay reinforcement, graph analysis, centrality metrics, elastic modulus, surface modification
\end{IEEEkeywords}

\section{Introduction}

Polymer-clay nanocomposites have emerged as a prominent class of advanced materials, offering significant improvements in mechanical, thermal, and barrier properties over conventional polymer systems \cite{ray2003polymer, azeez2013epoxy}. The incorporation of nanoscale clay platelets, particularly montmorillonite (MMT), into polymer matrices has been extensively investigated over the past two decades, yielding a vast body of experimental data on structure-property relationships \cite{okada2006twenty, ray2014recent}.

Despite the wealth of available research, systematic approaches for synthesizing this knowledge remain limited. Traditional literature reviews provide qualitative insights but often fail to capture the quantitative relationships between material parameters and property enhancements. Furthermore, the effectiveness of surface modification strategies---a critical factor in nanocomposite performance---varies considerably across different polymer systems, necessitating a more nuanced analytical framework \cite{chee2019effect, manias2001polypropylene}.

The mechanical reinforcement of polymers through clay incorporation is governed by several factors, including the aspect ratio and dispersion state of clay platelets, the interfacial adhesion between matrix and filler, and the intrinsic stiffness of the polymer matrix \cite{zare2017approach, zairi2011study}. The interphase region between clay and polymer plays a particularly crucial role, with its properties significantly influencing the overall composite behavior \cite{heydari2016multiscale, zare2020development}.

Graph-based analytical methods have proven effective in capturing complex relationships within large datasets across various scientific domains. In materials science, network representations can encode the relationships between processing conditions, material compositions, and resulting properties, enabling the identification of patterns that may not be apparent through traditional statistical analysis.

This study introduces a graph-theoretic framework for analyzing polymer-clay nanocomposite data through modulus-based clustering and bipartite network construction. The methodology enables systematic identification of polymer systems that serve as ``knowledge hubs''---materials with extensive experimental characterization demonstrating substantial property improvements. By employing centrality metrics borrowed from network science, we quantify the relative importance of different polymer matrices within their respective stiffness regimes.

\section{Background and Related Work}

\subsection{Mechanical Properties of Polymer-Clay Nanocomposites}

The mechanical reinforcement mechanisms in polymer-clay nanocomposites have been extensively studied through both experimental and computational approaches. Rao and Pochan \cite{rao2007mechanics} demonstrated that the mechanical properties of polymer-clay systems depend critically on the clay dispersion state, with exfoliated morphologies yielding superior reinforcement compared to intercalated structures. Similarly, Chan et al. \cite{chan2011mechanism} elucidated the reinforcement mechanism in nanoclay/polymer composites, highlighting the role of stress transfer at the clay-matrix interface.

Epoxy-clay nanocomposites have received particular attention due to epoxy's widespread use in structural applications. Akbari and Bagheri \cite{akbari2007deformation} investigated the deformation mechanisms in epoxy/clay systems, while Yasmin et al. \cite{yasmin2006mechanical} characterized the mechanical and thermal behavior of clay-reinforced epoxy. The synthesis parameters significantly influence final properties, as demonstrated by Fakhreddini-Najafabadi et al. \cite{fakhreddini2021investigation}, who systematically investigated the effects of synthesis conditions on mechanical performance.

\subsection{Surface Modification Effects}

Surface modification of clay nanoparticles represents a critical strategy for improving interfacial adhesion and overall composite properties. Organomodification of MMT using quaternary ammonium compounds enhances compatibility with hydrophobic polymer matrices and promotes intercalation or exfoliation \cite{sinha2005role}. Rauschendorfer et al. \cite{rauschendorfer2020tuning} demonstrated that surface-functionalized MMT nanosheets enable tuning of mechanical properties in poly(methyl acrylate) systems, achieving substantial modulus improvements through controlled interfacial chemistry.

The effectiveness of surface modification varies with polymer matrix characteristics. Harshita et al. \cite{harshita2018extrapolation} extrapolated mechanical strengthening effects in nanoclay/epoxy systems to elevated temperatures, while Su et al. \cite{su2024functionalized} showed that functionalized MMT/epoxy systems exhibit enhanced thermal and mechanical properties. For polyamide systems, Naveau et al. \cite{naveau2011effect} demonstrated that clay modification significantly influences the structure and mechanical properties of PA6 nanocomposites.

\subsection{Modeling Approaches}

Various modeling approaches have been employed to predict and understand the mechanical properties of polymer-clay nanocomposites. Zare and Rhee \cite{zare2020simulation} developed models for Young's modulus considering mechanical percolation and clay-interphase networks, while their subsequent work \cite{zare2022development} addressed halloysite nanotube nanocomposites through interphase zone analysis. Heydari-Meybodi et al. \cite{heydari2015new} proposed a multiscale approach for predicting elastic modulus that accounts for interfacial debonding.

Machine learning approaches have recently been applied to nanocomposite property prediction. Zazoum et al. \cite{zazoum2020modeling} employed deep neural networks for modeling mechanical properties of clay-reinforced nanocomposites, while Akbarzadeh et al. \cite{akbarzadeh2024machine} utilized machine learning for optimization of PLA/PBAT/clay blend nanocomposites.

\section{Methodology}

\subsection{Dataset Compilation and Clustering Basis}

The foundation of this analysis rests upon a curated dataset comprising 927 nanocomposite samples extracted from peer-reviewed literature sources. These samples were categorized based on surface modification status, yielding 789 samples (85.1\%) with modified nanofillers and 138 samples (14.9\%) with unmodified nanofillers. Prior to analysis, the dataset underwent preprocessing steps including removal of entries with incomplete mechanical property data, standardization of numeric formats, and handling of missing values through listwise deletion when interpolation was not scientifically justifiable.

The clustering strategy adopted in this work diverges from conventional unsupervised approaches. Rather than applying k-means or hierarchical clustering algorithms, we employed a domain-informed classification scheme based on the elastic modulus of the neat polymer matrix ($E_m$). This physically-motivated approach ensures that comparisons between nanocomposite systems occur within mechanically similar polymer categories, thereby controlling for baseline matrix stiffness when evaluating nanofiller reinforcement effects.

Four distinct clusters were established according to the following modulus ranges, as summarized in Table~\ref{tab:clusters}:

\begin{table}[htbp]
\caption{Modulus-Based Cluster Definitions}
\begin{center}
\begin{tabular}{@{}llcc@{}}
\toprule
Cluster & Material Type & Modulus Range (GPa) & Samples \\
\midrule
C1 & Soft elastomeric & $0 < E_m \leq 0.1$ & 23 \\
C2 & Semi-soft thermoplastic & $0.1 < E_m \leq 0.5$ & 47 \\
C3 & Intermediate & $0.5 < E_m \leq 1.0$ & 82 \\
C4 & Rigid polymer & $E_m > 1.0$ & 757 \\
\bottomrule
\end{tabular}
\label{tab:clusters}
\end{center}
\end{table}

Cluster C1 (Soft Elastomeric) encompasses natural rubber, silicone elastomers, and carboxylated nitrile rubber (CNBR) systems \cite{abhisha2020effect}. Cluster C2 (Semi-soft Thermoplastic) includes polycaprolactone (PCL), low-density polyethylene, and poly(methyl acrylate) variants \cite{kassa2020effects, carli2011phbv}. Cluster C3 (Intermediate) captures materials such as high-density polyethylene and polypropylene \cite{ebadi2015experimental, esmizadeh2018fabrication}. Cluster C4 (Rigid Polymer) constitutes the largest cluster with engineering thermoplastics including polyamides, epoxy resins, and poly(methyl methacrylate) \cite{wilkinson2007tensile, gabr2018mechanical, stojsic2016study}.

The pronounced imbalance toward Cluster C4 reflects the research community's emphasis on rigid matrix systems for structural applications, a pattern that itself provides insight into current research priorities in the nanocomposites field.

\subsection{Bipartite Graph Construction}

To capture the relationships between polymer matrices and their corresponding nanocomposite formulations, we constructed bipartite graphs for each modulus-based cluster. A bipartite graph $G = (U, V, E)$ consists of two disjoint vertex sets where edges connect only vertices from different sets, making this structure naturally suited to representing polymer-composite relationships.

The vertex sets are defined as follows:

\textbf{Polymer Matrix Nodes ($U$):} Each unique polymer matrix material within a cluster is represented as a node. These nodes encode the material identity and average elastic modulus within the cluster.

\textbf{Composite Sample Nodes ($V$):} Individual nanocomposite samples constitute the second node set. Each sample node encodes the specific combination of matrix, nanofiller type, loading fraction, processing conditions, and modification status reported in the source literature.

The edge set $E$ connects each composite sample to its corresponding base polymer matrix. For a composite sample $v_j$ fabricated from polymer matrix $u_i$, an edge $e_{ij} \in E$ is established with weight defined as:

\begin{equation}
w(e_{ij}) = \left| \frac{E_c - E_m}{E_m} \times 100 \right|
\label{eq:weight}
\end{equation}

where $E_c$ denotes the elastic modulus of the nanocomposite and $E_m$ represents the modulus of the neat polymer matrix. This weighting scheme quantifies the absolute percentage improvement in elastic modulus achieved through nanofiller incorporation, regardless of whether the change is positive or negative.

Visual distinction between modification states was implemented through node coloring in the generated graph visualizations. Composite samples prepared with surface-modified nanofillers are rendered in green, while those utilizing unmodified nanofillers appear in red.

\subsection{Centrality Evaluation Metrics}

Centrality measures provide quantitative assessment of node importance within graph structures. Three complementary centrality metrics were computed to characterize polymer matrix nodes and identify materials serving as focal points for nanocomposite research.

\textbf{Degree Centrality} measures the fraction of nodes to which a given node is connected. For a polymer matrix node $v$ in a graph with $n$ total nodes:

\begin{equation}
C_D(v) = \frac{\deg(v)}{n-1}
\label{eq:degree}
\end{equation}

where $\deg(v)$ represents the number of edges incident to node $v$. High degree centrality indicates that numerous nanocomposite formulations have been investigated using that matrix material.

\textbf{Weighted Degree Centrality} extends the basic degree measure by incorporating edge weights:

\begin{equation}
C_{WD}(v) = \sum_{u \in N(v)} w(v, u)
\label{eq:weighted_degree}
\end{equation}

This metric serves as the primary indicator for identifying ``knowledge hubs'' within each cluster. A polymer matrix exhibiting high weighted degree centrality has demonstrated substantial cumulative property improvements across its associated nanocomposite systems.

\textbf{Betweenness Centrality} quantifies the extent to which a node lies on shortest paths between other node pairs:

\begin{equation}
C_B(v) = \sum_{s \neq v \neq t} \frac{\sigma_{st}(v)}{\sigma_{st}}
\label{eq:betweenness}
\end{equation}

where $\sigma_{st}$ denotes the total number of shortest paths from node $s$ to node $t$, and $\sigma_{st}(v)$ represents the number of those paths passing through $v$.

\section{Results}

\subsection{Cluster Characteristics and Graph Statistics}

The four modulus-based clusters exhibit distinct characteristics in terms of sample distribution, polymer diversity, and network structure. Table~\ref{tab:graph_stats} summarizes the graph statistics for each cluster.

\begin{table}[htbp]
\caption{Graph Statistics by Cluster}
\begin{center}
\begin{tabular}{@{}lcccc@{}}
\toprule
Metric & C1 & C2 & C3 & C4 \\
\midrule
Polymer types & 4 & 9 & 12 & 49 \\
Total nodes & 27 & 56 & 94 & 806 \\
Total edges & 23 & 47 & 82 & 757 \\
Avg. improvement (\%) & 351.3 & 288.9 & 42.8 & 33.9 \\
\bottomrule
\end{tabular}
\label{tab:graph_stats}
\end{center}
\end{table}

A clear inverse relationship emerges between matrix stiffness and improvement magnitude. Soft elastomeric systems (C1) achieve an average improvement of 351.3\%, while rigid polymers (C4) show more modest improvements of 33.9\% on average. This pattern aligns with theoretical expectations: softer matrices have greater capacity for relative stiffness enhancement through the incorporation of rigid clay platelets.

Fig.~\ref{fig:cluster_graphs} presents the bipartite network visualizations for all four clusters, illustrating the polymer-composite relationships and modification patterns within each stiffness regime.

\begin{figure*}[htbp]
\centering
\begin{tabular}{cc}
\includegraphics[width=0.48\textwidth]{cluster_C1_graph.png} &
\includegraphics[width=0.48\textwidth]{cluster_C2_graph.png} \\
(a) Cluster C1: Soft elastomeric & (b) Cluster C2: Semi-soft thermoplastic \\[6pt]
\includegraphics[width=0.48\textwidth]{cluster_C3_graph.png} &
\includegraphics[width=0.48\textwidth]{cluster_C4_graph.png} \\
(c) Cluster C3: Intermediate & (d) Cluster C4: Rigid polymer \\
\end{tabular}
\caption{Bipartite network graphs for each modulus-based cluster. Blue nodes represent polymer matrices, green nodes indicate modified nanocomposite samples, and red nodes denote unmodified samples. Edge colors reflect improvement direction (green: positive, red: negative), with edge width proportional to improvement magnitude. The network density increases from C1 to C4, reflecting the greater research emphasis on rigid polymer systems.}
\label{fig:cluster_graphs}
\end{figure*}

Fig.~\ref{fig:cluster_summary} provides a complementary view showing the distribution of elastic modulus improvements versus matrix modulus for each cluster.

\begin{figure}[htbp]
\centerline{\includegraphics[width=\columnwidth]{cluster_summary.png}}
\caption{Scatter plots of elastic modulus improvement (\%) versus matrix modulus (GPa) for each cluster. Green markers represent modified samples; red markers represent unmodified samples. Note the different y-axis scales reflecting the substantially higher improvements achieved in softer polymer systems (C1, C2).}
\label{fig:cluster_summary}
\end{figure}

\subsection{Knowledge Hub Identification}

The centrality analysis identified distinct knowledge hubs within each cluster, as presented in Table~\ref{tab:centrality}.

\begin{table}[htbp]
\caption{Top Polymer Knowledge Hubs by Cluster}
\begin{center}
\begin{tabular}{@{}llccc@{}}
\toprule
Cluster & Polymer & $C_D$ & $C_{WD}$ & $C_B$ \\
\midrule
\multirow{3}{*}{C1} & PA6 & 0.154 & 3822.86 & 0.019 \\
 & CNBR & 0.154 & 2076.42 & 0.019 \\
 & NBR & 0.154 & 1648.19 & 0.019 \\
\midrule
\multirow{3}{*}{C2} & PMA & 0.109 & 12318.18 & 0.010 \\
 & PCL & 0.327 & 527.77 & 0.103 \\
 & Epoxy & 0.109 & 296.67 & 0.010 \\
\midrule
\multirow{3}{*}{C3} & Epoxy & 0.237 & 1517.09 & 0.054 \\
 & Nylon 6 & 0.043 & 511.25 & 0.001 \\
 & PP & 0.065 & 416.67 & 0.004 \\
\midrule
\multirow{3}{*}{C4} & Epoxy & 0.268 & 6386.10 & 0.072 \\
 & PA6 & 0.143 & 5284.98 & 0.020 \\
 & Nylon 6 & 0.068 & 2569.24 & 0.005 \\
\bottomrule
\end{tabular}
\label{tab:centrality}
\end{center}
\end{table}

In Cluster C1, PA6 (polyamide 6) emerges as the primary knowledge hub with the highest weighted degree centrality (3822.86), indicating exceptional property improvements in soft matrix formulations. This finding is particularly noteworthy as PA6 is typically classified as a rigid polymer; however, specific low-modulus PA6 grades or plasticized variants within the soft regime demonstrate remarkable reinforcement potential, consistent with findings by Abdel-Gawad et al. \cite{abdel2022fabrication}.

Cluster C2 identifies PMA (poly(methyl acrylate)) as the dominant knowledge hub with an exceptionally high weighted degree centrality of 12318.18. This value reflects extensive experimental characterization showing consistent, substantial property improvements in PMA-clay systems, as documented by Rauschendorfer et al. \cite{rauschendorfer2020tuning}.

For both intermediate (C3) and rigid (C4) polymer regimes, epoxy systems emerge as the primary knowledge hubs. Epoxy exhibits the highest degree centrality across these clusters (0.237 and 0.268, respectively), indicating extensive experimental investigation. The high weighted degree values (1517.09 and 6386.10) reflect both the breadth of research and the consistent property improvements achieved in epoxy-clay nanocomposites \cite{azeez2013epoxy, shettar2020experimental, tcherbi2016effects}.

\subsection{Surface Modification Effectiveness}

A key finding of this analysis concerns the stiffness-dependent effectiveness of surface modification strategies. Table~\ref{tab:modification} presents the comparison between modified and unmodified systems across clusters.

\begin{table}[htbp]
\caption{Modified vs. Unmodified System Comparison}
\begin{center}
\begin{tabular}{@{}lccccc@{}}
\toprule
Cluster & \multicolumn{2}{c}{Sample Count} & \multicolumn{2}{c}{Avg. Improvement (\%)} & Diff. \\
\cmidrule(r){2-3} \cmidrule(r){4-5}
 & Mod. & Unmod. & Mod. & Unmod. & (\%) \\
\midrule
C1 & 14 & 9 & 320.7 & 414.1 & -93.4 \\
C2 & 44 & 3 & 320.8 & 98.1 & +222.6 \\
C3 & 70 & 12 & 46.0 & 33.5 & +12.5 \\
C4 & 643 & 114 & 31.5 & 42.1 & -10.6 \\
\bottomrule
\end{tabular}
\label{tab:modification}
\end{center}
\end{table}

Surface modification provides the most significant benefit in Cluster C2 (semi-soft thermoplastics), where modified systems achieve 222.6\% higher improvement than unmodified counterparts. This substantial difference suggests that interfacial adhesion plays a critical role in reinforcement efficiency for this stiffness regime.

Conversely, in both the softest (C1) and stiffest (C4) regimes, unmodified systems show higher average improvements. For C1, this may reflect the inherent compatibility of some elastomeric matrices with unmodified clays or the specific experimental conditions employed. For C4, the rigid polymer matrices may provide sufficient constraint on clay platelets without requiring additional surface functionalization, consistent with observations by Merah and Mohamed \cite{merah2019nanoclay}.

\section{Discussion}

The network-based analysis reveals several important patterns in polymer-clay nanocomposite behavior that have implications for both fundamental understanding and practical material design.

\subsection{Stiffness-Dependent Reinforcement Efficiency}

The inverse relationship between matrix stiffness and relative improvement magnitude can be understood through composite mechanics principles. For a polymer with initial modulus $E_m$, the addition of rigid clay platelets with modulus $E_f \gg E_m$ creates a composite whose modulus is bounded by the Voigt and Reuss limits. The relative improvement achievable is inherently greater when the initial matrix modulus is low, as the modulus contrast between filler and matrix is larger \cite{heydari2016multiscale, zare2017efficiency}.

This finding suggests that researchers seeking maximum relative property enhancement should focus on softer polymer systems, while those requiring high absolute stiffness values should employ rigid matrices where the improvements, though proportionally smaller, still yield substantial absolute gains.

\subsection{Knowledge Hub Implications}

The identification of knowledge hubs through centrality analysis provides guidance for future research directions. Polymer systems with high weighted degree centrality represent well-characterized material platforms with demonstrated reinforcement capability. For applications requiring reliable property enhancement, these systems offer reduced uncertainty compared to less-studied alternatives.

However, the analysis also reveals opportunities in underexplored regions. Polymer matrices with low centrality values but residing in high-potential clusters (such as C1 and C2) represent targets for future investigation. The exceptional performance of PMA in C2, for instance, suggests that other semi-soft thermoplastics may yield similar benefits upon systematic exploration.

\subsection{Modification Strategy Selection}

The stiffness-dependent effectiveness of surface modification has practical implications for nanocomposite design. For semi-soft thermoplastic applications, investment in surface modification chemistry is strongly justified by the substantial performance gains. For rigid polymer systems, the benefits of modification are more modest, and cost-benefit considerations may favor simpler processing approaches using unmodified clays.

This pattern may reflect the changing relative importance of different reinforcement mechanisms across stiffness regimes. In softer matrices, interfacial adhesion (enhanced by surface modification) may dominate the reinforcement mechanism, while in rigid matrices, the geometric constraint provided by well-dispersed clay platelets may be sufficient regardless of interfacial chemistry \cite{zare2020effects, domun2015improving}.

\section{Conclusion}

This study demonstrated the utility of graph-theoretic analysis for systematically evaluating polymer-clay nanocomposite research data. The modulus-based clustering approach provides a physically meaningful framework for comparing materials within similar stiffness regimes, while bipartite graph construction captures the relationships between polymer matrices and their composite formulations.

Key findings include:
\begin{itemize}
\item An inverse relationship between matrix stiffness and improvement magnitude, with soft systems (C1-C2) achieving 288-351\% average improvements compared to 33-46\% for rigid polymers (C3-C4)
\item Identification of knowledge hubs including PA6 (C1), PMA (C2), and Epoxy (C3-C4) as polymer systems with substantial cumulative property improvements
\item Stiffness-dependent surface modification effectiveness, with semi-soft thermoplastics (C2) showing the greatest benefit from modification (+222.6\%)
\end{itemize}

The methodology presented here provides a systematic framework for guiding material selection and identifying research opportunities in the nanocomposites field. Future work will extend this approach to include additional property dimensions and incorporate predictive modeling capabilities.

\begin{thebibliography}{00}
\bibitem{ray2003polymer} S. S. Ray and M. Okamoto, ``Polymer/layered silicate nanocomposites: A review from preparation to processing,'' \textit{Prog. Polym. Sci.}, vol. 28, no. 11, pp. 1539--1641, Nov. 2003.
\bibitem{azeez2013epoxy} A. A. Azeez, K. Y. Rhee, S. J. Park, and D. Hui, ``Epoxy clay nanocomposites - Processing, properties and applications: A review,'' \textit{Compos. Part B: Eng.}, vol. 45, no. 1, pp. 308--320, 2013.
\bibitem{okada2006twenty} A. Okada and A. Usuki, ``Twenty years of polymer-clay nanocomposites,'' \textit{Macromol. Mater. Eng.}, vol. 291, no. 12, pp. 1449--1476, Dec. 2006.
\bibitem{ray2014recent} S. S. Ray, ``Recent trends and future outlooks in the field of clay-containing polymer nanocomposites,'' \textit{Macromol. Chem. Phys.}, vol. 215, no. 12, pp. 1162--1179, 2014.
\bibitem{chee2019effect} S. S. Chee and M. Jawaid, ``The effect of bi-functionalized MMT on morphology, thermal stability, dynamic mechanical, and tensile properties of epoxy/organoclay nanocomposites,'' \textit{Polymers}, vol. 11, no. 12, Dec. 2019.
\bibitem{manias2001polypropylene} E. Manias, A. Touny, L. Wu, K. Strawhecker, B. Lu, and T. C. Chung, ``Polypropylene/montmorillonite nanocomposites. Review of the synthetic routes and materials properties,'' \textit{Chem. Mater.}, vol. 13, no. 10, pp. 3516--3523, 2001.
\bibitem{zare2017approach} Y. Zare, ``An approach to study the roles of percolation threshold and interphase in tensile modulus of polymer/clay nanocomposites,'' \textit{J. Colloid Interface Sci.}, vol. 486, pp. 249--254, Jan. 2017.
\bibitem{zairi2011study} F. Za\"{i}ri, J. M. Gloaguen, M. Na\"{i}t-Abdelaziz, A. Mesbah, and J. M. Lefebvre, ``Study of the effect of size and clay structural parameters on the yield and post-yield response of polymer/clay nanocomposites via a multiscale micromechanical modelling,'' \textit{Acta Mater.}, vol. 59, no. 10, pp. 3851--3863, 2011.
\bibitem{heydari2016multiscale} M. Heydari-Meybodi, S. Saber-Samandari, and M. Sadighi, ``3D multiscale modeling to predict the elastic modulus of polymer/nanoclay composites considering realistic interphase property,'' \textit{Compos. Interfaces}, vol. 23, no. 7, pp. 641--661, Sep. 2016.
\bibitem{zare2020development} Y. Zare and K. Y. Rhee, ``Development of a model for modulus of polymer halloysite nanotube nanocomposites by the interphase zones around dispersed and networked nanotubes,'' \textit{Sci. Rep.}, vol. 12, no. 1, Dec. 2022.
\bibitem{rao2007mechanics} Y. Q. Rao and J. M. Pochan, ``Mechanics of polymer-clay nanocomposites,'' \textit{Macromolecules}, vol. 40, no. 2, pp. 290--296, Jan. 2007.
\bibitem{chan2011mechanism} M. L. Chan, K. T. Lau, T. T. Wong, M. P. Ho, and D. Hui, ``Mechanism of reinforcement in a nanoclay/polymer composite,'' \textit{Compos. Part B: Eng.}, vol. 42, no. 6, pp. 1708--1712, Sep. 2011.
\bibitem{akbari2007deformation} B. Akbari and R. Bagheri, ``Deformation mechanism of epoxy/clay nanocomposite,'' \textit{Eur. Polym. J.}, vol. 43, no. 3, pp. 782--788, Mar. 2007.
\bibitem{yasmin2006mechanical} A. Yasmin, J. J. Luo, J. L. Abot, and I. M. Daniel, ``Mechanical and thermal behavior of clay/epoxy nanocomposites,'' \textit{Compos. Sci. Technol.}, vol. 66, no. 14, pp. 2415--2422, Nov. 2006.
\bibitem{fakhreddini2021investigation} S. Fakhreddini-Najafabadi, M. Torabi, and F. Taheri-Behrooz, ``An investigation on the effects of synthesis on the mechanical properties of nanoclay/epoxy,'' \textit{J. Mater. Res. Technol.}, vol. 15, pp. 5375--5395, Nov. 2021.
\bibitem{sinha2005role} S. Sinha Ray, K. Yamada, M. Okamoto, and K. Ueda, ``New polylactide/layered silicate nanocomposites,'' \textit{Polymer}, vol. 44, pp. 857--866, 2003.
\bibitem{rauschendorfer2020tuning} J. E. Rauschendorfer, K. M. Thien, M. Denz, S. K\"{o}ster, F. Ehlers, and P. Vana, ``Tuning the mechanical properties of poly(methyl acrylate) via surface-functionalized montmorillonite nanosheets,'' \textit{Macromol. Mater. Eng.}, vol. 2000595, pp. 1--9, Feb. 2020.
\bibitem{harshita2018extrapolation} P. N. Harshita, D. K. Rathore, R. K. Prusty, and B. C. Ray, ``Extrapolation of mechanical strengthening effect in nanoclay/epoxy nanocomposites to elevated temperature environments,'' \textit{Trans. Indian Inst. Met.}, vol. 71, no. 8, pp. 2015--2024, Aug. 2018.
\bibitem{su2024functionalized} L. Su, C. Fang, and H. Luo, ``Functionalized montmorillonite/epoxy resin nanocomposites with enhanced thermal and mechanical properties,'' \textit{RSC Adv.}, vol. 14, no. 42, pp. 31251--31258, Oct. 2024.
\bibitem{naveau2011effect} E. Naveau et al., ``Effect of clay modification on the structure and mechanical properties of polyamide-6 nanocomposites,'' \textit{Eur. Polym. J.}, vol. 47, no. 1, pp. 5--15, Jan. 2011.
\bibitem{zare2020simulation} Y. Zare and K. Y. Rhee, ``Simulation of Young's modulus for clay-reinforced nanocomposites assuming mechanical percolation, clay-interphase networks and interfacial linkage,'' \textit{J. Mater. Res. Technol.}, vol. 9, no. 6, pp. 12473--12483, Nov. 2020.
\bibitem{zare2022development} Y. Zare and K. Y. Rhee, ``Development of a model for modulus of polymer halloysite nanotube nanocomposites by the interphase zones around dispersed and networked nanotubes,'' \textit{Sci. Rep.}, vol. 12, no. 1, Dec. 2022.
\bibitem{heydari2015new} M. Heydari-Meybodi, S. Saber-Samandari, and M. Sadighi, ``A new approach for prediction of elastic modulus of polymer/nanoclay composites by considering interfacial debonding: Experimental and numerical investigations,'' \textit{Compos. Sci. Technol.}, vol. 117, pp. 379--385, Sep. 2015.
\bibitem{zazoum2020modeling} B. Zazoum, E. Triki, and A. Bachri, ``Modeling of mechanical properties of clay-reinforced polymer nanocomposites using deep neural network,'' \textit{Materials}, vol. 13, no. 19, Oct. 2020.
\bibitem{akbarzadeh2024machine} M. Akbarzadeh et al., ``Machine learning-driven optimization and prediction of poly(lactic acid)/poly(butylene-adipate-co-terephthalate)/clay blend nanocomposites in structure–property relationships,'' \textit{Polym. Eng. Sci.}, Mar. 2024.
\bibitem{abhisha2020effect} V. S. Abhisha, A. Augustine, J. Joseph, S. P. Thomas, and R. Stephen, ``Effect of halloysite nanotubes and organically modified bentonite clay hybrid filler system on the properties of natural rubber,'' \textit{J. Elastomers Plast.}, vol. 52, no. 5, pp. 432--452, Aug. 2020.
\bibitem{kassa2020effects} A. Kassa, A. Benhamida, M. Kaci, and S. Bruzaud, ``Effects of montmorillonite, sepiolite, and halloysite clays on the morphology and properties of polycaprolactone bionanocomposites,'' \textit{Polym. Polym. Compos.}, vol. 28, no. 5, pp. 338--347, Jun. 2020.
\bibitem{carli2011phbv} L. N. Carli, J. S. Crespo, and R. S. Mauler, ``PHBV nanocomposites based on organomodified montmorillonite and halloysite: The effect of clay type on the morphology and thermal and mechanical properties,'' \textit{Compos. Part A: Appl. Sci. Manuf.}, vol. 42, no. 11, pp. 1601--1608, Nov. 2011.
\bibitem{ebadi2015experimental} H. Ebadi-Dehaghani, H. A. Khonakdar, M. Barikani, and S. H. Jafari, ``Experimental and theoretical analyses of mechanical properties of PP/PLA/clay nanocomposites,'' \textit{Compos. Part B: Eng.}, vol. 69, pp. 133--144, 2015.
\bibitem{esmizadeh2018fabrication} E. Esmizadeh, G. Naderi, S. Arezoomand, and S. Mazinani, ``Fabrication and characterization properties of polypropylene/polycarbonate/clay nanocomposites prepared with twin-screw extruder,'' \textit{Sci. Eng. Compos. Mater.}, vol. 25, no. 1, pp. 31--39, Jan. 2018.
\bibitem{wilkinson2007tensile} A. N. Wilkinson et al., ``Tensile properties of melt intercalated polyamide 6 - Montmorillonite nanocomposites,'' \textit{Compos. Sci. Technol.}, vol. 67, no. 15--16, pp. 3360--3368, Dec. 2007.
\bibitem{gabr2018mechanical} M. H. Gabr and K. Uzawa, ``Mechanical, thermal, and interfacial shear properties of polyamide/nanoclay nanocomposites: Study the effect of nanoclay on the micromechanism of fracture,'' \textit{J. Thermoplast. Compos. Mater.}, vol. 31, no. 4, pp. 447--464, Apr. 2018.
\bibitem{stojsic2016study} J. Stoj\v{s}i\'{c}, P. Raos, and A. Kalendova, ``A study of structure and tensile properties of polyamide 12/clay nanocomposites,'' \textit{Polym. Compos.}, vol. 37, no. 3, pp. 684--691, Mar. 2016.
\bibitem{shettar2020experimental} M. Shettar, C. S. S. Kowshik, M. Manjunath, and P. Hiremath, ``Experimental investigation on mechanical and wear properties of nanoclay-epoxy composites,'' \textit{J. Mater. Res. Technol.}, vol. 9, no. 4, pp. 9108--9116, Jan. 2020.
\bibitem{tcherbi2016effects} A. Tcherbi-Narteh, M. Hosur, and S. Jeelani, ``Effects of different montmorillonite nanoclay loading on cure behavior and properties of diglycidyl ether of bisphenol A epoxy,'' \textit{J. Nanomater.}, vol. 2016, 2016.
\bibitem{abdel2022fabrication} A. M. Abdel-Gawad, A. R. Ramadan, A. Flores, and A. M. K. Esawi, ``Fabrication of Nylon 6-Montmorillonite clay nanocomposites with enhanced structural and mechanical properties by solution compounding,'' \textit{Polymers}, vol. 14, no. 21, Nov. 2022.
\bibitem{zare2017efficiency} Y. Zare, M. Fasihi, and K. Y. Rhee, ``Efficiency of stress transfer between polymer matrix and nanoplatelets in clay/polymer nanocomposites,'' \textit{Appl. Clay Sci.}, vol. 143, pp. 265--272, Jul. 2017.
\bibitem{merah2019nanoclay} N. Merah and O. Mohamed, ``Nanoclay and water uptake effects on mechanical properties of unsaturated polyester,'' \textit{J. Nanomater.}, vol. 2019, 2019.
\bibitem{zare2020effects} Y. Zare and K. Y. Rhee, ``Effects of critical interfacial shear modulus between polymer matrix and nanoclay on the effective interphase properties and tensile modulus of nanocomposites,'' \textit{Constr. Build. Mater.}, vol. 247, Jun. 2020.
\bibitem{domun2015improving} N. Domun, H. Hadavinia, T. Zhang, T. Sainsbury, G. H. Liaghat, and S. Vahid, ``Improving the fracture toughness and the strength of epoxy using nanomaterials-a review of the current status,'' \textit{Nanoscale}, vol. 7, no. 23, pp. 10294--10329, Jun. 2015.
\end{thebibliography}

\end{document}
